% concatenated from A_complete_classification_of_endomorphisms_of_Kiselmans_semigroup.tex
%Version 3.1 December 2024
% See section 11 of the User Manual for version history
%
%%%%%%%%%%%%%%%%%%%%%%%%%%%%%%%%%%%%%%%%%%%%%%%%%%%%%%%%%%%%%%%%%%%%%%
%%                                                                 %%
%% Please do not use \input{...} to include other tex files.       %%
%% Submit your LaTeX manuscript as one .tex document.              %%
%%                                                                 %%
%% All additional figures and files should be attached             %%
%% separately and not embedded in the \TeX\ document itself.       %%
%%                                                                 %%
%%%%%%%%%%%%%%%%%%%%%%%%%%%%%%%%%%%%%%%%%%%%%%%%%%%%%%%%%%%%%%%%%%%%%

%%\documentclass[referee,sn-basic]{sn-jnl}% referee option is meant for double line spacing

%%=======================================================%%
%% to print line numbers in the margin use lineno option %%
%%=======================================================%%

%%\documentclass[lineno,pdflatex,sn-basic]{sn-jnl}% Basic Springer Nature Reference Style/Chemistry Reference Style

%%=========================================================================================%%
%% the documentclass is set to pdflatex as default. You can delete it if not appropriate.  %%
%%=========================================================================================%%

%%\documentclass[sn-basic]{sn-jnl}% Basic Springer Nature Reference Style/Chemistry Reference Style

%%Note: the following reference styles support Namedate and Numbered referencing. By default the style follows the most common style. To switch between the options you can add or remove “Numbered” in the optional parenthesis. 
%%The option is available for: sn-basic.bst, sn-chicago.bst%  
 
%%\documentclass[pdflatex,sn-nature]{sn-jnl}% Style for submissions to Nature Portfolio journals
%%\documentclass[pdflatex,sn-basic]{sn-jnl}% Basic Springer Nature Reference Style/Chemistry Reference Style
\documentclass[pdflatex,sn-mathphys-num]{sn-jnl}% Math and Physical Sciences Numbered Reference Style
%%\documentclass[pdflatex,sn-mathphys-ay]{sn-jnl}% Math and Physical Sciences Author Year Reference Style
%%\documentclass[pdflatex,sn-aps]{sn-jnl}% American Physical Society (APS) Reference Style
%%\documentclass[pdflatex,sn-vancouver-num]{sn-jnl}% Vancouver Numbered Reference Style
%%\documentclass[pdflatex,sn-vancouver-ay]{sn-jnl}% Vancouver Author Year Reference Style
%%\documentclass[pdflatex,sn-apa]{sn-jnl}% APA Reference Style
%%\documentclass[pdflatex,sn-chicago]{sn-jnl}% Chicago-based Humanities Reference Style

%%%% Standard Packages
%%<additional latex packages if required can be included here>

\usepackage{graphicx}%
\usepackage[T1]{fontenc}
\usepackage{multirow}%
\usepackage{amsmath,amssymb,amsfonts}%
\usepackage{amsthm}%
\usepackage{mathrsfs}%
\usepackage[title]{appendix}%
\usepackage{xcolor}%
\usepackage{textcomp}%
\usepackage{manyfoot}%
\usepackage{booktabs}%
\usepackage{algorithm}%
\usepackage{algorithmicx}%
\usepackage{algpseudocode}%
\usepackage{listings}%
%%%%

%%%%%=============================================================================%%%%
%%%%  Remarks: This template is provided to aid authors with the preparation
%%%%  of original research articles intended for submission to journals published 
%%%%  by Springer Nature. The guidance has been prepared in partnership with 
%%%%  production teams to conform to Springer Nature technical requirements. 
%%%%  Editorial and presentation requirements differ among journal portfolios and 
%%%%  research disciplines. You may find sections in this template are irrelevant 
%%%%  to your work and are empowered to omit any such section if allowed by the 
%%%%  journal you intend to submit to. The submission guidelines and policies 
%%%%  of the journal take precedence. A detailed User Manual is available in the 
%%%%  template package for technical guidance.
%%%%%=============================================================================%%%%

%% as per the requirement new theorem styles can be included as shown below
\theoremstyle{thmstyleone}%
\newtheorem{theorem}{Theorem}%  meant for continuous numbers
%%\newtheorem{theorem}{Theorem}[section]% meant for sectionwise numbers
%% optional argument [theorem] produces theorem numbering sequence instead of independent numbers for Proposition
\newtheorem{proposition}[theorem]{Proposition}% 
%%\newtheorem{proposition}{Proposition}% to get separate numbers for theorem and proposition etc.

\theoremstyle{thmstyletwo}%
\newtheorem{example}{Example}%
\newtheorem{remark}{Remark}%

\theoremstyle{thmstylethree}%
\newtheorem{definition}{Definition}%

\newtheorem{corollary}[theorem]{Corollary}

\raggedbottom
%%\unnumbered% uncomment this for unnumbered level heads

\begin{document}

\title[A complete classification of endomorphisms of Kiselman's semigroup]{A complete classification of endomorphisms of Kiselman's semigroup}

%%=============================================================%%
%% GivenName	-> \fnm{Joergen W.}
%% Particle	-> \spfx{van der} -> surname prefix
%% FamilyName	-> \sur{Ploeg}
%% Suffix	-> \sfx{IV}
%% \author*[1,2]{\fnm{Joergen W.} \spfx{van der} \sur{Ploeg} 
%%  \sfx{IV}}\email{iauthor@gmail.com}
%%=============================================================%%

\author*[1]{\fnm{Luka} \sur{Andren\v sek}}\email{andrensek.luka@gmail.com}





\affil*[1]{\orgdiv{Department of Mathematics}, \orgname{Faculty of Mathematics and Physics, University of Ljubljana}, \orgaddress{\street{Jadranska ulica 19}, \city{Ljubljana}, \postcode{SI-1000},  \country{Slovenia}}}



%%==================================%%
%% Sample for unstructured abstract %%
%%==================================%%

\abstract{Kiselman's semigroup $K_n$ was studied by Kudryavtseva and Mazorchuk, who posed the question of whether it is possible to classify all endomorphisms of $K_n$. In this paper, we provide a complete classification of endomorphisms of $K_n$ and present a Boolean matrix monoid that is isomorphic to $\text{End}(K_n)$.}

%%================================%%
%% Sample for structured abstract %%
%%================================%%

% \abstract{\textbf{Purpose:} The abstract serves both as a general introduction to the topic and as a brief, non-technical summary of the main results and their implications. The abstract must not include subheadings (unless expressly permitted in the journal's Instructions to Authors), equations or citations. As a guide the abstract should not exceed 200 words. Most journals do not set a hard limit however authors are advised to check the author instructions for the journal they are submitting to.
% 
% \textbf{Methods:} The abstract serves both as a general introduction to the topic and as a brief, non-technical summary of the main results and their implications. The abstract must not include subheadings (unless expressly permitted in the journal's Instructions to Authors), equations or citations. As a guide the abstract should not exceed 200 words. Most journals do not set a hard limit however authors are advised to check the author instructions for the journal they are submitting to.
% 
% \textbf{Results:} The abstract serves both as a general introduction to the topic and as a brief, non-technical summary of the main results and their implications. The abstract must not include subheadings (unless expressly permitted in the journal's Instructions to Authors), equations or citations. As a guide the abstract should not exceed 200 words. Most journals do not set a hard limit however authors are advised to check the author instructions for the journal they are submitting to.
% 
% \textbf{Conclusion:} The abstract serves both as a general introduction to the topic and as a brief, non-technical summary of the main results and their implications. The abstract must not include subheadings (unless expressly permitted in the journal's Instructions to Authors), equations or citations. As a guide the abstract should not exceed 200 words. Most journals do not set a hard limit however authors are advised to check the author instructions for the journal they are submitting to.}

\keywords{Kiselman's semigroup, endomorphism, monoid, semigroup}

%%\pacs[JEL Classification]{D8, H51}

%%\pacs[MSC Classification]{35A01, 65L10, 65L12, 65L20, 65L70}

\maketitle

\section{Introduction and the main result}\label{sec1}

Kiselman's semigroup was first introduced by Kiselman in the context of convexity theory in \cite{Kiselman}. A generalisation was later introduced and thoroughly studied by 
Kudryavtseva and Mazorchuk in \cite{GK}. The semigroup arises naturally in various settings such as graph dynamics as discussed by Collina and D’Andrea in \cite{Graphs}.

In \cite{ganyushkin11}, Ganyushkin and Mazorchuk introduced Kiselman quotients of 0-Hecke monoids associated with simply laced Dynkin diagrams, and
more generally, Hecke--Kiselman monoids associated with finite simple digraphs. Further properties and identities of Hecke--Kiselman monoids were studied by Ashikhmin, Volkov and Zhang in \cite{Ashikhmin15},  
irreducible representations of Hecke--Kiselman monoids were studied by Wiertel in \cite{wiertel22},
and the word problem for some particular Hecke--Kiselman monoids was studied by Lebed in \cite{Lebed2025}. 
Furthermore, the finiteness of Hecke--Kiselman monoids was investigated by Aragona and D’Andrea in 
\cite{aragona13}, while D’Andrea and Stella showed in \cite{dandrea23} that the cardinality of Kiselman's semigroups grows double-exponentially. 
 
Kiselman's semigroup is also related to 0-Hecke algebras \cite{norton79}. Properties and structure of Hecke--Kiselman algebras were studied by various authors in \cite{okninski20, okninski21, wiertel23, Mecel2019}.  

%, their radicals in \cite{okninski21}, their Gelfand-Kirillov dimension was discussed in \cite{wiertel23}, and their Gröbner bases and automaton properties were discussed in \cite{mecel19}.
In order to formulate the main result of our paper, we define the following notions.
Let $M \in \mathbb{R}^{m \times n}$ be a real $m \times n$ matrix. Then for $1 \leq i \leq m$ and $1 \leq j \leq n$ we denote by $M_{i, j}$ the entry of the matrix $M$ in the $i$-th row and the $j$-th column. 
\begin{definition}
   Let $A \in \mathbb{R}^{2 \times 2}$ and $M \in \mathbb{R}^{n \times n}$ be real matrices with $n \geq 2$. We say that $A$ is a \emph{$2 \times 2$ submatrix} of 
   $M$ if there exist columns $i, j \in \{1, 2, \dots, n\}$ and rows $x, y \in \{1, 2, \dots, n\}$ with $i < j$ and $x < y$ such that 
   \[
      A = 
      \begin{pmatrix}
         M_{x, i} & M_{x, j} \\
         M_{y, i} & M_{y, j}
      \end{pmatrix}.
   \]
\end{definition}

Now we define the set of Boolean matrices that avoid a certain submatrix.

\begin{definition}\label{Dn}
   Let $n \geq 2$ be an integer, and let
   \begin{equation}\label{pattern}
      P = 
      \begin{pmatrix}
         0 & 1 \\
         1 & 0
      \end{pmatrix}.
   \end{equation}
   We define the set $D_n$ as 
   \[
      D_n = \{ M \in \{0, 1\}^{n \times n} \mid P \text{ is not a } 2 \times 2 \text{ submatrix of } M \}.
   \]
\end{definition}

On the set $D_n$, we define a binary operation $\cdot$ as follows. Let $A, B \in D_n$. Then $C = A \cdot B$, with
\begin{align*}
   C_{i, j} = \bigvee_{k=1}^n (A_{i, k} \land B_{k, j}).
\end{align*}
Here, $\lor$ denotes the Boolean join and $\land$ denotes the Boolean meet on $\{0, 1\}$. One can easily verify that $(D_n, \cdot)$ is a monoid, and its unit element is the identity matrix. 

Let $\mathcal{A} = \{a_1, a_2, \dots , a_n\}$ be a finite alphabet. Define Kiselman's semigroup $K_n$ by
\begin{align*}
   K_n = \langle a_1, a_2, \dots, a_n \mid a_i^2 = a_i, \; a_i a_j a_i = a_j a_i a_j = a_j a_i, \; 1 \leq i < j \leq n \rangle .
\end{align*}

In \cite{GK}, the authors showed in Proposition 20 that the only automorphism of $K_n$ is the identity, and that the only antiautomorphism is the unique extension to an antiautomorphism of the map 
$a_i \mapsto a_{n-i+1}$. They leave open the question (Question~21) concerning the classification of endomorphisms of $K_n$, which has not yet been answered.

The main result of this paper is the following theorem.
\begin{theorem}\label{thm:main}
   Let $n \geq 2$ be an integer. Then $\operatorname{End}(K_n) \cong (D_n, \cdot)$.
\end{theorem}

\section{Proof of the main result}
\subsection{Properties of idempotents of $K_n$}

Let $X \subseteq \{1, 2, \dots , n\}$. If $X = \emptyset$ define $e_X = e$, the unit element in $K_n$. Otherwise, write $X = \{i_1, i_2, \dots , i_k \mid i_1 > i_2 > \dots > i_k\}$. Then set
$e_X = a_{i_1} a_{i_2} \dots a_{i_k}$. Now define $E_n = \{e_X \mid X \subseteq \{1, 2, \dots, n\}\}$. In Proposition 11 of \cite{GK}, it is shown that every element of $E_n$ is an idempotent, that the elements of $E_n$ are pairwise distinct, and that every idempotent of $K_n$ is of the form $e_X$ for some $X \subseteq \{1, 2, \dots , n\}$. Consequently, the number of idempotents in $K_n$ is precisely $2^n$. 


In \cite[Proposition 14]{GK}, the following proposition was proved. 

\begin{proposition}\label{idem:eq}
   Let $X, Y \subseteq \{1, 2, \dots, n\}$. Then the following conditions are equivalent:
   \begin{enumerate}
      \item $e_X e_Y$ is an idempotent.
      \item $e_X e_Y = e_{X \cup Y}$.
      \item For every $x \in X \setminus Y$ and every $y \in Y \setminus X$ we have $x > y$.
   \end{enumerate}
\end{proposition}


Define the \emph{content} map $c : K_n \rightarrow 2^{\{1, 2, \dots, n\}}$, where $c(a)$ is the set of all $i$'s such that $a_i$ appears in $a$. In \cite{GK}, it is shown that this map is well defined.
Furthermore, in Lemma 10 of \cite{GK} it is shown that the following holds.


\begin{proposition}
  The map $c$ is a semigroup epimorphism from $K_n$ onto the semigroup $(2^{\{1, 2, \dots, n\}}, \cup)$.
\end{proposition} 

We now prove the following proposition.
\begin{proposition}\label{gate}
   Let $X, Y \subseteq \{1, 2, \dots , n\}$. The following two conditions are equivalent:
   \begin{enumerate}
      \item For every $x \in X \setminus Y$ and every $y \in Y \setminus X$ we have $x > y$.
      \item $e_X e_Y$ is an idempotent and $e_X e_Y e_X = e_Y e_X e_Y = e_X e_Y$.
   \end{enumerate}
\end{proposition} 


\begin{proof}
   The implication \((2) \Rightarrow (1)\) is immediate from Proposition~\ref{idem:eq}.

   By Proposition~\ref{idem:eq}, to prove \((1) \Rightarrow (2)\), it suffices to show the equalities:
   \[
   e_X e_Y e_X = e_Y e_X e_Y = e_X e_Y.
   \]

   Since $e_X e_Y$ is an idempotent, it is equal to $e_{X \cup Y}$. Observe that condition $(3)$ in Proposition~\ref{idem:eq} holds for the pairs of sets $(X \cup Y, X)$ and $(Y, X \cup Y)$, respectively. It follows that 
   \begin{align*}
      e_X e_Y e_X = (e_X e_Y) e_X = e_{X \cup Y} e_X = e_{(X \cup Y) \cup X} = e_{X \cup Y} = e_X e_Y,
   \end{align*}
   and similarly
   \begin{align*}
      e_Y e_X e_Y = e_Y (e_X e_Y ) = e_Y e_{X \cup Y } = e_{Y \cup (X \cup Y)} = e_{X \cup Y} = e_X e_Y.
   \end{align*}
   This concludes the proof.
\end{proof}

\subsection{Isomorphism of $\operatorname{End}(K_n)$ with the monoid of monotone sequences of sets $M_n$}

In this section, we construct an isomorphism from $\text{End}(K_n)$ to the set of sequences of sets satisfying the \textit{monotonicity} condition, which we now define. 

\begin{definition}
   Let $\mathbf{X} = (X_1, X_2, \dots, X_n)$ be a sequence of sets, where $X_i \subseteq \{1, 2, \dots, n\}$ for $1 \leq i \leq n$. We say that $\mathbf{X}$ is \emph{monotone} if
   \begin{align*}
      &\text{for all } i, j \in \{1, 2, \dots, n\} \text{ with } j > i, \\
      &\text{and for all } x \in X_j \setminus X_i \text{ and } y \in X_i \setminus X_j, \text{ we have } x > y.
   \end{align*}
   We denote by $M_n$ the set
   \[
   M_n = \{(X_1, X_2, \dots, X_n) \mid (X_1, X_2, \dots, X_n) \text{ is monotone}\}.
   \]
\end{definition}



On the set $M_n$, we define a binary operation $\ast$ as follows.
Let $\mathbf{X} = (X_1, X_2, \dots, X_n)$ and $\mathbf{Y} = (Y_1, Y_2, \dots, Y_n)$ be elements of $M_n$. 
We define $\mathbf{Z} = \mathbf{X} \ast \mathbf{Y}$, where $\mathbf{Z} = (Z_1, Z_2, \dots, Z_n)$ is given by
\begin{align*}
   Z_i = \bigcup\limits_{j \in Y_i} X_j
\end{align*}
for $1 \leq i \leq n$. It is easy to verify that $(M_n, \ast)$ is a monoid, with its unit element given by the sequence $(\{1\}, \{2\}, \dots,\{n\})$.



\begin{theorem}\label{thm:Phi}
   The map $\Phi : \operatorname{End}(K_n) \rightarrow M_n$, defined by
   \begin{align*}
      \Phi(\varphi) = \big( c(\varphi(a_1)), c(\varphi(a_2)), \dots, c(\varphi(a_n)) \big)
   \end{align*}
   is an isomorphism of monoids $\operatorname{End}(K_n)$ and $(M_n, \ast)$.
\end{theorem}



\begin{proof}
   First, we show that for every $\varphi \in \operatorname{End}(K_n)$, $\Phi(\varphi)$ lies in $M_n$. Let $\varphi \in \operatorname{End}(K_n)$. 
   Since $\varphi$ is a homomorphism, it maps idempotents to idempotents. In particular, $\varphi(a_i)$ is an idempotent for 
   $1 \leq i \leq n$. As every idempotent in $K_n$ is of the form $e_X$, it follows that there exists a sequence $(X_1, X_2, \dots, X_n)$ of subsets of $\{1, 2, \dots, n\}$ such that
   \begin{align*}
      \varphi(a_i) = e_{X_i}
   \end{align*}
   for $1 \leq i \leq n$. Thus we have $\Phi(\varphi) = (X_1, X_2, \dots, X_n)$. Now observe that $a_j a_i$ is an idempotent, whenever $j > i$, and in this case we also have $a_ja_ia_j = a_i a_j a_i = a_j a_i$. Applying $\varphi$ to these equalities and using the fact that $\varphi$ preserves idempotents yields:
   \begin{align*}
      & e_{X_j} e_{X_i} \text{ is an idempotent and} \\
      & e_{X_j} e_{X_i} e_{X_j} = e_{X_i} e_{X_j} e_{X_i} = e_{X_j} e_{X_i}
   \end{align*}
   whenever $j > i$. By Proposition~\ref{gate} it follows that for every $j>i$, the condition 
   \begin{align*}
      \text{for every } x \in X_j \setminus X_i\; \text{and every } y \in X_i \setminus X_j\; \text{we have } x > y
   \end{align*}
   must be satisfied. This implies the monotonicity of $(X_1, X_2, \dots, X_n)$ and hence $\Phi(\varphi) \in M_n$. 





   The fact that $\Phi$ is injective follows quickly, since for any endomorphism of $K_n$, the images of the generators are of the form $e_{X_i}$ for some $X_i \subseteq \{1, 2, \dots, n\}$.
   Observe that an idempotent in $K_n$ is uniquely determined by its content. Since an endomorphism is uniquely determined by the images of the generators, it follows that an endomorphism of $K_n$ is uniquely determined by the contents of the images of the generators, from which the injectivity of $\Phi$ follows.


   Now we prove surjectivity. Let $\mathbf{X} = (X_1, X_2, \dots, X_n) \in M_n$. Define the map $\psi : \{a_1, a_2, \dots, a_n\} \rightarrow K_n$ by 
   \begin{align*}
      \psi(a_i) = e_{X_i}.
   \end{align*}
   Set $b_i = \psi(a_i)$ for $1 \leq i \leq n$. Clearly, each $b_i$ is an idempotent. Furthermore, let $j > i$. Then, by Proposition~\ref{gate} and the monotonicity condition imposed on $(X_1, X_2, \dots, X_n)$ it must follow that
   \begin{align*}
      b_j b_i b_j = b_i b_j b_i = b_j b_i.
   \end{align*}
   Thus, the elements $\{b_1, b_2, \dots, b_n\}$ satisfy the defining relations of $K_n$, so we can uniquely extend $\psi$ to an endomorphism of $K_n$. By construction, we have that $\Phi(\psi) = \mathbf{X}$, which proves surjectivity. 

%~

   Now we prove that $\Phi$ is a homomorphism. Let $\psi, \varphi \in \operatorname{End}(K_n)$. Let
   $\psi(a_j) = e_{Y_j}$ and $\varphi(a_j) = e_{X_j}$ for $1 \leq j \leq n$. Also write $e_{X_i} = a_{m_l} a_{m_{l-1}} \dots a_{m_1}$, where $X_i = \{m_l, m_{l-1}, \dots, m_1\}$ and $m_l > m_{l-1} > \dots > m_1$. We compute
   \begin{align*}
      \hspace{-1.3cm} (\psi \circ \varphi)(a_i) &= \psi(e_{X_i}) \\
      &= \psi(a_{m_l} a_{m_{l-1}} \dots a_{m_2} a_{m_1})  \\
      &= \psi(a_{m_l}) \psi(a_{m_{l-1}}) \dots \psi(a_{m_2}) \psi(a_{m_1}) \\
      &= e_{Y_{m_l}} e_{Y_{m_{l-1}}} \dots e_{Y_{m_2}} e_{Y_{m_1}} \\
      &= e_{\cup_{1 \leq k \leq l} Y_{m_k}}.
   \end{align*}
   The last equality holds since $(\psi \circ \varphi)(a_i)$ is an idempotent, every idempotent in $K_n$ is of the form $e_X$, and the content of an element in $K_n$ is preserved under its defining relations, since the content map $c$ is well defined. Now, if we denote by
   $(\Phi(\psi \circ \varphi))_i$ the $i$-th component of the sequence of sets $\Phi(\psi \circ \varphi)$, we get
   
   \begin{align*}
      \hspace{-6cm} \big(\Phi(\psi \circ \varphi)\big)_i &= c((\psi \circ \varphi)(a_i)) \\
      &= c(e_{\cup_{1 \leq k \leq l} Y_{m_k}}) \\
      &= \bigcup\limits_{k =1}^{l} Y_{m_k} \\
      &= \bigcup\limits_{j \in X_i} Y_{j},
   \end{align*}
   since $X_i = \{m_l, m_{l-1}, \dots, m_1\}$.

   On the other hand, one has
   \begin{align*}
      \big(\Phi(\psi) \ast \Phi(\varphi)\big)_i &= \big((c(\psi(a_1)), c(\psi(a_2)), \dots, c(\psi(a_n))) \ast (c(\varphi(a_1)), c(\varphi(a_2)), \dots, c(\varphi(a_n)))\big)_i \\
      &= \big((Y_1, Y_2, \dots, Y_n) \ast (X_1, X_2, \dots, X_n)\big)_i \\
      &= \bigcup\limits_{j \in X_i} Y_j.
   \end{align*}
   We have shown that $\big(\Phi(\psi \circ \varphi)\big)_i = \big(\Phi(\psi) \ast \Phi(\varphi)\big)_i$,
   and since $i$ was chosen arbitrarily, it follows that
   \begin{align*}
      \Phi(\psi \circ \varphi) = \Phi(\psi) \ast \Phi(\varphi).
   \end{align*}
   Since the identity map is the unit of $\operatorname{End}(K_n)$, and because of
   \begin{align*}
      \Phi(\operatorname{id}_{K_n}) = (\{1\}, \{2\}, \dots, \{n\}),
   \end{align*}
   we conclude that $\Phi$ is an isomorphism of monoids. This completes the proof. 

\end{proof}

\subsection{Isomorphism of $M_n$ with $D_n$}

One can conclude that an $n \times n$ Boolean matrix $M$ is an element of $D_n$, as defined in Definition~\ref{Dn},  for $n \geq 2$, if and only if there do not exist columns 
$i, j \in \{1, 2, \dots, n\}$ and rows $x, y \in \{1, 2, \dots, n\}$ with $i < j$, and $x < y$, such that
\begin{align*}
   &M_{x, i} = 0, \quad M_{x, j} = 1, \\
   &M_{y, i} = 1, \quad M_{y, j} = 0.
\end{align*}
Also observe that if $A$ and $B$ are $n \times n$ Boolean matrices, and if for all $i, j \in \{1, 2, \dots, n\}$ we have that 
$A_{i, j} = 1$ if and only if $B_{i, j} = 1$, then $A = B$.  

Let $(e_1, e_2, \dots, e_n)$ be the standard orthonormal basis of $\mathbb{R}^n$. Define the map 
$\chi : 2^{\{1, 2, \dots, n\}} \rightarrow \{0, 1\}^n$ by
\begin{align*}
   \chi(X) = \sum_{i \in X} e_i.
\end{align*}
If further $v \in \mathbb{R}^n$ is a vector, then for $1 \leq i \leq n$ we denote by $v_i$ the $i$-th element of the vector $v$. 

\begin{theorem}\label{thm:Psi}
   The map $\Psi : M_n \rightarrow D_n$, defined by
   \[
      \Psi(X_1, X_2, \dots, X_n) = (\chi(X_1), \chi(X_2), \dots, \chi(X_n))
   \]
   where the right-hand side is interpreted as a matrix whose $i$-th column is $\chi(X_i)$, is an isomorphism of monoids $(M_n, \ast)$ and $(D_n, \cdot)$.
\end{theorem}

\begin{proof}
   We first show that $\Psi$ maps into $D_n$. Let $\mathbf{X} = (X_1, X_2, \dots, X_n) \in M_n$ and denote $M = \Psi(\mathbf{X})$.
   Observe that $M_{x, i} = 1$ if and only if $(\chi(X_i))_x = 1$, which is equivalent to $x \in X_i$.
   Now assume $i < j$ and assume that the following equalities hold
   \begin{align*}
      &M_{x, i} = 0, \quad M_{x, j} = 1, \\
      &M_{y, i} = 1, \quad M_{y, j} = 0.
   \end{align*}
   We aim to show that $x < y$ is impossible. 
   The upper equalities are equivalent to 
   \begin{align*}
      &x \notin X_i, \quad x \in X_j, \\
      &y \in X_i, \quad y \notin X_j, 
   \end{align*}
   which is further equivalent to 
   \[ 
      x \in X_j \setminus X_i \quad \text{and} \quad y \in X_i \setminus X_j. 
   \]
   We apply monotonicity of $(X_1, X_2, \dots, X_n)$ and get $x>y$. Hence, $\Psi(\mathbf{X}) \in D_n$.



   The map $\Psi$ is obviously injective.

   We now prove surjectivity. Let $M \in D_n$. Define $X_i = \{x \in \{1, 2, \dots, n\} \mid M_{x, i} = 1\}$. First, we check that $(X_1, X_2, \dots, X_n)$ is monotone.
   Let $i < j$ and suppose $x \in X_j \setminus X_i$ and $y \in X_i \setminus X_j$. Then we have
   \begin{align*}
      M_{x, i} = 0, \quad M_{x, j} = 1, \\
      M_{y, i} = 1, \quad M_{y, j} = 0.
   \end{align*}
   Since $M \in D_n$, it must follow that $x>y$, and hence $(X_1, X_2, \dots, X_n)$ is monotone.
   Furthermore, for $x, i \in \{1, 2, \dots, n\}$ we have that $(\Psi(X_1, X_2, \dots, X_n))_{x, i} = 1$ if and only if $x \in X_i$, which is equivalent to  $M_{x, i} = 1$ by the definition of $X_i$.
   Hence
   \[
   \Psi(X_1, X_2, \dots, X_n) = M,
   \]
   and surjectivity is proven. 


   Now we show that $\Psi$ is a homomorphism. Let $\mathbf{X}, \mathbf{Y} \in M_n$, where
   $\mathbf{X} = (X_1, X_2, \dots, X_n)$ and $\mathbf{Y} = (Y_1, Y_2, \dots, Y_n)$. We compute
   \begin{align*}
      \Psi(\mathbf{X} \ast \mathbf{Y}) = \Psi(\bigcup_{j \in Y_1} X_j, \bigcup_{j \in Y_2} X_j, \dots, \bigcup_{j \in Y_n} X_j).
   \end{align*}
   It follows that for $x, i \in \{1, 2, \dots, n\}$ we have 
   \begin{align*}
      \Psi(\mathbf{X} \ast \mathbf{Y})_{x, i} = 1 &\Leftrightarrow x \in \bigcup_{j\in Y_i}X_j \\
      &\Leftrightarrow \text{there exists } k \in Y_i \text{ such that } x \in X_k.
   \end{align*}
   On the other hand,
   \begin{align*}
      (\Psi(\mathbf{X}) \cdot \Psi(\mathbf{Y}))_{x, i} = \bigvee_{k=1}^n (\Psi(\mathbf{X})_{x, k} \land \Psi(\mathbf{Y})_{k, i}).
   \end{align*}
   Hence, we have
   \begin{align*}
      (\Psi(\mathbf{X}) \cdot \Psi(\mathbf{Y}))_{x, i} = 1 &\Leftrightarrow \text{there exists } k \in \{1, 2, \dots, n\} \text{ such that } \Psi(\mathbf{X})_{x, k} = \Psi(\mathbf{Y})_{k, i} = 1 \\
      &\Leftrightarrow \text{there exists } k \in \{1, 2, \dots, n\} \text{ such that } x \in X_k \text{ and } k \in Y_i \\
      &\Leftrightarrow \text{there exists } k \in Y_i \text{ such that } x \in X_k.
   \end{align*}
   We thus obtain
   \[
      \Psi(\mathbf{X} \ast \mathbf{Y})_{x, i} = (\Psi(\mathbf{X}) \cdot \Psi(\mathbf{Y}))_{x, i}.
   \]
   Since $x$ and $i$ were chosen arbitrarily, we conclude that
   \begin{align*}
      \Psi(\mathbf{X} \ast \mathbf{Y}) = \Psi(\mathbf{X}) \cdot \Psi(\mathbf{Y}).
   \end{align*}
   Furthermore, since $\Psi(\{1\}, \{2\}, \dots, \{n\}) = I$, it follows that $\Psi$ is an isomorphism of monoids. This concludes the proof.

\end{proof}

Now we are in a position to prove the main result of the paper.

\begin{proof}[Proof of Theorem~\ref{thm:main}]
   Theorem~\ref{thm:Phi} gives the isomorphism
   \[
      \Phi : \operatorname{End}(K_n) \rightarrow M_n.
   \]
   Theorem~\ref{thm:Psi} shows that
   \[
      \Psi : M_n \rightarrow D_n
   \]
   is also an isomorphism.  Composing these two maps yields an isomorphism
   \[
      \Psi \circ \Phi : \operatorname{End}(K_n) \rightarrow D_n,
   \]
   so \(\operatorname{End}(K_n) \cong (D_n,\cdot)\).  This completes the proof of the main theorem.  \qedhere
\end{proof}



\section{Concluding remarks}

Even though the following proposition was proved in \cite{GK} in Proposition 20, we can prove it again with relative ease.
\begin{proposition}
   Let $n\in\mathbb{N}$. The only invertible element of $\operatorname{End}(K_n)$ is the identity. Consequently, $\operatorname{Aut}(K_n)$ is trivial.
\end{proposition}

\begin{proof}
   In \cite{Rutherford} it is shown that the only invertible $n \times n$ Boolean matrices are the permutation matrices.
   Since the only permutation matrix that is an element of $D_n$ is the identity matrix, it follows that the only invertible element of $D_n$ is the identity matrix. This concludes the proof.
\end{proof}

Finding an explicit formula for the cardinality of $\operatorname{End}(K_n)$ or equivalently of $D_n$ has proven to be challenging. The authors in \cite{Kitaev} study matrices avoiding certain patterns. If we denote by $c_{m, n}$ the number of $m \times n$ Boolean matrices avoiding the pattern (submatrix) \eqref{pattern}, they proposed a finite algorithm which, for each fixed $m \geq 1$, gives a closed formula for $c_{m, n}$ for any $n \geq 1$. This result (Corollary 13 in \cite{Kitaev}) is summarised in the following proposition.
\begin{proposition}
   Let $n \in \mathbb{N}$. Then we have
   \begin{itemize}
      \item $c_{2, n} = (3 + n) \cdot 3^{n-1}$,
      \vspace{0.2cm}
      \item $c_{3, n} = \frac{1}{3} \cdot (2 + n)(96 + 31n + n^2) \cdot 4^{n - 3}$,
      \vspace{0.2cm}
      \item $c_{4, n} = \frac{1}{36} \cdot (2812500 + 3963450n + 1862971n^2 + 339300n^3 + 21265n^4 + 510n^5 + 4n^6) \cdot 5^{n-7}$.
   \end{itemize}
\end{proposition} 
From this, one can determine the cardinality of $D_n$ for $2 \leq n \leq 4$. 






\backmatter


\bmhead{Acknowledgements}

I would like to thank Ganna Kudryavtseva for her valuable mentorship, insightful guidance, and helpful suggestions throughout the preparation of this paper.
I would also like to thank the anonymous referee for their careful reading and valuable suggestions.

%%===========================================================================================%%
%
%\nocite{sn-bib-options}
%\bibliography{citations1}% common bib file

%% BioMed_Central_Bib_Style_v1.01

\begin{thebibliography}{17}
% BibTex style file: bmc-mathphys.bst (version 2.1), 2014-07-24
\ifx \bisbn   \undefined \def \bisbn  #1{ISBN #1}\fi
\ifx \binits  \undefined \def \binits#1{#1}\fi
\ifx \bauthor  \undefined \def \bauthor#1{#1}\fi
\ifx \batitle  \undefined \def \batitle#1{#1}\fi
\ifx \bjtitle  \undefined \def \bjtitle#1{#1}\fi
\ifx \bvolume  \undefined \def \bvolume#1{\textbf{#1}}\fi
\ifx \byear  \undefined \def \byear#1{#1}\fi
\ifx \bissue  \undefined \def \bissue#1{#1}\fi
\ifx \bfpage  \undefined \def \bfpage#1{#1}\fi
\ifx \blpage  \undefined \def \blpage #1{#1}\fi
\ifx \burl  \undefined \def \burl#1{\textsf{#1}}\fi
\ifx \doiurl  \undefined \def \doiurl#1{\url{https://doi.org/#1}}\fi
\ifx \betal  \undefined \def \betal{\textit{et al.}}\fi
\ifx \binstitute  \undefined \def \binstitute#1{#1}\fi
\ifx \binstitutionaled  \undefined \def \binstitutionaled#1{#1}\fi
\ifx \bctitle  \undefined \def \bctitle#1{#1}\fi
\ifx \beditor  \undefined \def \beditor#1{#1}\fi
\ifx \bpublisher  \undefined \def \bpublisher#1{#1}\fi
\ifx \bbtitle  \undefined \def \bbtitle#1{#1}\fi
\ifx \bedition  \undefined \def \bedition#1{#1}\fi
\ifx \bseriesno  \undefined \def \bseriesno#1{#1}\fi
\ifx \blocation  \undefined \def \blocation#1{#1}\fi
\ifx \bsertitle  \undefined \def \bsertitle#1{#1}\fi
\ifx \bsnm \undefined \def \bsnm#1{#1}\fi
\ifx \bsuffix \undefined \def \bsuffix#1{#1}\fi
\ifx \bparticle \undefined \def \bparticle#1{#1}\fi
\ifx \barticle \undefined \def \barticle#1{#1}\fi
\bibcommenthead
\ifx \bconfdate \undefined \def \bconfdate #1{#1}\fi
\ifx \botherref \undefined \def \botherref #1{#1}\fi
\ifx \url \undefined \def \url#1{\textsf{#1}}\fi
\ifx \bchapter \undefined \def \bchapter#1{#1}\fi
\ifx \bbook \undefined \def \bbook#1{#1}\fi
\ifx \bcomment \undefined \def \bcomment#1{#1}\fi
\ifx \oauthor \undefined \def \oauthor#1{#1}\fi
\ifx \citeauthoryear \undefined \def \citeauthoryear#1{#1}\fi
\ifx \endbibitem  \undefined \def \endbibitem {}\fi
\ifx \bconflocation  \undefined \def \bconflocation#1{#1}\fi
\ifx \arxivurl  \undefined \def \arxivurl#1{\textsf{#1}}\fi
\csname PreBibitemsHook\endcsname

%%% 1
\bibitem[\protect\citeauthoryear{Aragona and D’Andrea}{2013}]{aragona13}
\begin{barticle}
\bauthor{\bsnm{Aragona}, \binits{R.}},
\bauthor{\bsnm{D’Andrea}, \binits{A.}}:
\batitle{{Hecke-Kiselman} monoids of small cardinality}.
\bjtitle{Semigr. Forum}
\bvolume{86}(\bissue{1}),
\bfpage{32}--\blpage{40}
(\byear{2013})
\end{barticle}
\endbibitem

%%% 2
\bibitem[\protect\citeauthoryear{Ashikhmin et~al.}{2015}]{Ashikhmin15}
\begin{barticle}
\bauthor{\bsnm{Ashikhmin}, \binits{D.}},
\bauthor{\bsnm{Volkov}, \binits{M.}},
\bauthor{\bsnm{Zhang}, \binits{W.}}:
\batitle{The {Finite} {Basis} {Problem} for {Kiselman} {Monoids}}.
\bjtitle{Demonstr. Math.}
\bvolume{48}(\bissue{4}),
\bfpage{475}--\blpage{492}
(\byear{2015})
\end{barticle}
\endbibitem

%%% 3
\bibitem[\protect\citeauthoryear{Collina and D’Andrea}{2015}]{Graphs}
\begin{barticle}
\bauthor{\bsnm{Collina}, \binits{E.}},
\bauthor{\bsnm{D’Andrea}, \binits{A.}}:
\batitle{A graph-dynamical interpretation of {Kiselman’s} semigroups}.
\bjtitle{J. Algebr. Comb.}
\bvolume{41}(\bissue{4}),
\bfpage{1115}--\blpage{1132}
(\byear{2015})
\end{barticle}
\endbibitem

%%% 4
\bibitem[\protect\citeauthoryear{D'Andrea and Stella}{2023}]{dandrea23}
\begin{barticle}
\bauthor{\bsnm{D'Andrea}, \binits{A.}},
\bauthor{\bsnm{Stella}, \binits{S.}}:
\batitle{The cardinality of {Kiselman's} semigroups grows double-exponentially}.
\bjtitle{Bull. Belg. Math. Soc. Simon Stevin}
\bvolume{30}(\bissue{5}),
\bfpage{570}--\blpage{576}
(\byear{2023})
\end{barticle}
\endbibitem

%%% 5
\bibitem[\protect\citeauthoryear{Ganyushkin and Mazorchuk}{2011}]{ganyushkin11}
\begin{barticle}
\bauthor{\bsnm{Ganyushkin}, \binits{O.}},
\bauthor{\bsnm{Mazorchuk}, \binits{V.}}:
\batitle{On {Kiselman} quotients of 0-{Hecke} monoids}.
\bjtitle{Int. Electron. J. Algebra}
\bvolume{10}(\bissue{10}),
\bfpage{174}--\blpage{191}
(\byear{2011})
\end{barticle}
\endbibitem

%%% 6
\bibitem[\protect\citeauthoryear{Kiselman}{2002}]{Kiselman}
\begin{barticle}
\bauthor{\bsnm{Kiselman}, \binits{C.}}:
\batitle{A semigroup of operators in convexity theory}.
\bjtitle{Trans. Am. Math. Soc.}
\bvolume{354}(\bissue{5}),
\bfpage{2035}--\blpage{2053}
(\byear{2002})
\end{barticle}
\endbibitem

%%% 7
\bibitem[\protect\citeauthoryear{Kitaev et~al.}{2005}]{Kitaev}
\begin{barticle}
\bauthor{\bsnm{Kitaev}, \binits{S.}},
\bauthor{\bsnm{Mansour}, \binits{T.}},
\bauthor{\bsnm{Vella}, \binits{A.}}:
\batitle{Pattern avoidance in matrices}.
\bjtitle{J. Integer Seq.}
\bvolume{8}(\bissue{2}),
\bfpage{05}--\blpage{2216052216}
(\byear{2005})
\end{barticle}
\endbibitem


%%% 8
\bibitem[\protect\citeauthoryear{Kudryavtseva and Mazorchuk}{2009}]{GK}
\begin{barticle}
\bauthor{\bsnm{Kudryavtseva}, \binits{G.}},
\bauthor{\bsnm{Mazorchuk}, \binits{V.}}:
\batitle{On {Kiselman’s} semigroup}.
\bjtitle{Yokohama Math. J.}
\bvolume{55}(\bissue{1}),
\bfpage{22}--\blpage{46}
(\byear{2009})
\end{barticle}
\endbibitem

%%% 9
\bibitem[\protect\citeauthoryear{Lebed}{2025}]{Lebed2025}
\begin{barticle}
\bauthor{\bsnm{Lebed}, \binits{V.}}:
\batitle{The word problem for {Hecke--Kiselman} monoids of type {$A_n$} and {$\widetilde A_n$}}.
\bjtitle{Semigr. Forum}
\bvolume{110}(\bissue{3}),
\bfpage{597}--\blpage{614}
(\byear{2025})
\end{barticle}
\endbibitem

%%% 10
\bibitem[\protect\citeauthoryear{M{\k{e}}cel and Okni{\'{n}}ski}{2019}]{Mecel2019}
\begin{barticle}
\bauthor{\bsnm{M{\k{e}}cel}, \binits{A.}},
\bauthor{\bsnm{Okni{\'{n}}ski}, \binits{J.}}:
\batitle{Gr{\"o}bner basis and the automaton property of {Hecke--Kiselman} algebras}.
\bjtitle{Semigr. Forum}
\bvolume{99}(\bissue{2}),
\bfpage{447}--\blpage{464}
(\byear{2019})
\end{barticle}
\endbibitem

%%% 11
\bibitem[\protect\citeauthoryear{Norton}{1979}]{norton79}
\begin{barticle}
\bauthor{\bsnm{Norton}, \binits{P.}}:
\batitle{0-{Hecke} algebras}.
\bjtitle{J. Aust. Math. Soc.}
\bvolume{27}(\bissue{3}),
\bfpage{337}--\blpage{357}
(\byear{1979})
\end{barticle}
\endbibitem

%%% 12
\bibitem[\protect\citeauthoryear{Okni{\'n}ski and Wiertel}{2020}]{okninski20}
\begin{barticle}
\bauthor{\bsnm{Okni{\'n}ski}, \binits{J.}},
\bauthor{\bsnm{Wiertel}, \binits{M.}}:
\batitle{Combinatorics and structure of {Hecke--Kiselman} algebras}.
\bjtitle{Commun. Contemp. Math.}
\bvolume{22}(\bissue{07}),
\bfpage{2050022}
(\byear{2020})
\end{barticle}
\endbibitem

%%% 13
\bibitem[\protect\citeauthoryear{Okni{\'n}ski and Wiertel}{2021}]{okninski21}
\begin{barticle}
\bauthor{\bsnm{Okni{\'n}ski}, \binits{J.}},
\bauthor{\bsnm{Wiertel}, \binits{M.}}:
\batitle{On the radical of a {Hecke--Kiselman} algebra}.
\bjtitle{Algebras Represent. Theory}
\bvolume{24}(\bissue{6}),
\bfpage{1431}--\blpage{1440}
(\byear{2021})
\end{barticle}
\endbibitem

%%% 14
\bibitem[\protect\citeauthoryear{Rutherford}{1963}]{Rutherford}
\begin{barticle}
\bauthor{\bsnm{Rutherford}, \binits{D.E.}}:
\batitle{Inverses of {Boolean} matrices}.
\bjtitle{Glasg. Math. J.}
\bvolume{6}(\bissue{1}),
\bfpage{49}--\blpage{53}
(\byear{1963})
\end{barticle}
\endbibitem

%%% 15
\bibitem[\protect\citeauthoryear{Wiertel}{2022}]{wiertel22}
\begin{barticle}
\bauthor{\bsnm{Wiertel}, \binits{M.}}:
\batitle{Irreducible representations of {Hecke--Kiselman} monoids}.
\bjtitle{Linear Algebra Its Appl.}
\bvolume{640},
\bfpage{12}--\blpage{33}
(\byear{2022})
\end{barticle}
\endbibitem

%%% 16
\bibitem[\protect\citeauthoryear{Wiertel}{2023}]{wiertel23}
\begin{barticle}
\bauthor{\bsnm{Wiertel}, \binits{M.}}:
\batitle{The {Gelfand--Kirillov} dimension of {Hecke--Kiselman} algebras}.
\bjtitle{Forum Math.}
\bvolume{35}(\bissue{2}),
\bfpage{523}--\blpage{534}
(\byear{2023})
\end{barticle}
\endbibitem

\end{thebibliography}

\end{document}
